% ============================================================================
% Fichier  : partie2/P02_C03_acquerir_preserver.tex
% Titre    : Acquérir et Préserver
% Auteur   : MINKA MI NGUIDJOÏ Thierry Emmanuel
% Version  : 1.0 -- novembre 2026
% Corrigés : corriges/Corr_C03.tex
% ============================================================================
% Sources : RFC 3227 (ordre de volatilité) ; ISO/IEC 27037 (identification,
% recueil, acquisition, préservation, paraphrasée) ; NIST SP 800-86.
% Outils cités (existants) : bloqueurs d'écriture, formats brut, E01 (EWF),
% AFF4 ; outils de capture mémoire WinPmem, DumpIt, LiME, AVML.
% Faits nouveaux (bible) : acquisition de SRV-CPT01 le 15/03, 16 h 00 à
% 19 h 25, E01, NC-03 ; SRV-FIC01 chiffré BitLocker, acquisition logique à
% chaud le 18/03 (NC-07) ; journaux de la passerelle : rotation à 60 jours.
% ============================================================================

\chapter{Acquérir et Préserver}
\label{chap:acquisition}
\chaptermark{Acquérir et préserver}

\epigraph{%
    « On n'examine jamais l'original.\\
    On examine une copie dont on peut prouver qu'elle lui est égale. »%
}{Principe du cours}

\niveauChapitre{Fondamental}{%
    Maîtrise du Chapitre~\ref{chap:preuve}
}

\begin{boxobjectifs}
\begin{itemize}
    \item Appliquer l'ordre de volatilité à une situation concrète.
    \item Choisir entre acquisition à froid et acquisition à chaud, et
          justifier ce choix.
    \item Acquérir un disque avec un bloqueur d'écriture, dans un format
          adapté, et vérifier l'acquisition.
    \item Capturer la mémoire vive en connaissant l'empreinte que laisse
          l'outil.
    \item Préserver des journaux avant leur rotation et des données
          détenues par des tiers.
    \item Documenter une acquisition de façon qu'un autre puisse la
          contrôler.
\end{itemize}
\end{boxobjectifs}

\bigskip

% ============================================================================
\section*{Avant de commencer : un serveur chiffré}
% ============================================================================

Bafoussam, mardi 18 mars 2025. Mme~MBARGA veut l'image du serveur de
fichiers SRV-FIC01, celui qui a envoyé 3,2~Go vers l'extérieur. Le serveur
est toujours allumé, débranché du réseau depuis le vendredi. M.~KAMGA
rappelle un détail~: son disque de données est chiffré par BitLocker, et la
clé de récupération est conservée chez le prestataire qui a installé la
machine, à Douala. Si l'on éteint le serveur pour en extraire le disque, on
obtiendra une copie parfaite d'un volume chiffré, illisible tant que la clé
n'arrive pas. Si on le laisse allumé, le volume est déverrouillé et lisible,
mais toute copie faite à chaud modifie un peu la machine.

\begin{boxalert}
\textbf{Violation typique.} Appliquer mécaniquement la règle « on éteint
puis on copie le disque » à une machine dont le chiffrement rendra le disque
illisible une fois éteinte.

\smallskip
\textbf{Règle~:} on choisit la méthode d'acquisition d'après ce que l'on
perdrait avec chacune, et l'on écrit la raison du choix. La meilleure
méthode est celle qui conserve le plus de preuve utile au prix de la plus
petite altération documentée.
\end{boxalert}

\separateur

% ============================================================================
\section{L'Ordre de Volatilité}
\label{sec:c03-volatilite}
% ============================================================================

La RFC~3227~\cite{rfc3227} recommande de recueillir les traces de la plus
volatile à la moins volatile. On en tire l'ordre suivant pour un serveur.

\begin{tableaucours}{cL{4.6cm}L{2.6cm}L{4.2cm}}{%
    \tableauHeader{Rang} &
    \tableauHeader{Trace} &
    \tableauHeader{Durée de vie} &
    \tableauHeader{Moyen d'acquisition}
}
1 & Registres, cache du processeur & Nanosecondes & Hors de portée en
pratique \\
2 & Mémoire vive~: processus, connexions, tables de routage et ARP, clés
en mémoire & Jusqu'à l'extinction & Capture de la mémoire \\
3 & Fichiers temporaires, fichier d'échange & Heures à jours & Image du
disque \\
4 & Disques & Mois à années & Image du disque \\
5 & Journaux distants, données de supervision & Selon la rotation &
Export, réquisition \\
6 & Configuration physique, topologie & Stable & Photographies, schéma \\
7 & Supports d'archive & Années & Copie \\
\end{tableaucours}
\vspace{-0.8em}
\begin{center}{\small\itshape Tableau~3.1. L'ordre de volatilité (d'après la RFC~3227)}\end{center}

Le rang~5 surprend souvent~: un journal distant paraît durable, alors que
sa rotation peut l'effacer en quelques jours. Les journaux de la passerelle
de NYANGA CRÉDIT ne gardent que 60~jours~: le premier accès du 2~mars
disparaîtrait le 1\textsuperscript{er}~mai. Leur export, le 14~mars, était
une urgence au même titre que la capture mémoire.

% ============================================================================
\section{À Froid ou à Chaud}
\label{sec:c03-froid-chaud}
% ============================================================================

\begin{boxdef}
\textbf{\termecle{Acquisition à froid}} (\textit{dead acquisition})~: copie
d'un support après extinction de la machine, à l'aide d'un bloqueur
d'écriture. Elle ne modifie pas le support et perd tout ce qui était en
mémoire.

\smallskip
\textbf{\termecle{Acquisition à chaud}} (\textit{live acquisition})~:
recueil sur une machine allumée (mémoire, volume déverrouillé, état du
système). Elle conserve le volatil et laisse une empreinte de l'outil sur la
machine, qu'il faut documenter.
\end{boxdef}

\begin{tableaucours}{L{4.4cm}L{3.8cm}L{4.0cm}}{%
    \tableauHeader{Situation} &
    \tableauHeader{Méthode} &
    \tableauHeader{Raison}
}
Serveur allumé, attaque possiblement en cours & À chaud d'abord (mémoire),
puis à froid & Processus et connexions perdus à l'extinction \\
Volume chiffré, clé indisponible & À chaud (copie logique du volume
déverrouillé) & Volume illisible une fois éteint \\
Poste éteint à l'arrivée & À froid & Rien de volatil à préserver \\
Machine virtuelle & Instantané de l'hyperviseur, puis copie des fichiers
de disque & Mémoire et disque figés ensemble, sans outil dans la machine \\
Service en nuage & Export par l'interface du fournisseur, réquisition &
Aucun accès au matériel \\
Serveur critique qui ne peut s'arrêter & À chaud, copie logique ciblée &
Arrêt impossible~; périmètre écrit et justifié \\
\end{tableaucours}
\vspace{-0.8em}
\begin{center}{\small\itshape Tableau~3.2. Choisir la méthode d'acquisition}\end{center}

% ============================================================================
\section{Acquérir un Disque}
\label{sec:c03-disque}
% ============================================================================

\subsection{Le bloqueur d'écriture}

Un système d'exploitation écrit sur tout support qu'on lui présente~: il met
à jour des dates d'accès, monte des volumes, écrit des journaux. Le
\termecle{bloqueur d'écriture} s'interpose entre le support et la station
d'acquisition et laisse passer les lectures en refusant les écritures. Le
bloqueur matériel est un boîtier placé sur le câble~; le bloqueur logiciel
est un réglage du système de la station. Le premier est préféré en enquête
judiciaire, parce que son comportement se teste et se documente
indépendamment du système. On vérifie un bloqueur avant usage, sur un
support de test, et l'on note la vérification.

\subsection{Les formats d'image}

\begin{tableaucours}{L{2.2cm}L{4.8cm}L{5.2cm}}{%
    \tableauHeader{Format} &
    \tableauHeader{Caractéristiques} &
    \tableauHeader{Usage}
}
Brut (\texttt{dd}, \texttt{.raw}) & Copie secteur par secteur, sans
métadonnées ni compression & Universel, lisible par tout outil~; empreinte
et notes dans des fichiers séparés \\
E01 (EWF) & Métadonnées intégrées (enquêteur, numéro de pièce, notes),
compression, contrôle par blocs, empreinte stockée & Format courant des
outils d'enquête~; bon compromis \\
AFF4 & Format ouvert, métadonnées riches, adapté aux très gros volumes et à
la mémoire & Moins répandu, utile pour les acquisitions volumineuses \\
\end{tableaucours}
\vspace{-0.8em}
\begin{center}{\small\itshape Tableau~3.3. Formats d'image}\end{center}

\subsection{Les pièges du stockage moderne}

Les disques à mémoire flash (SSD) gèrent eux-mêmes l'effacement des blocs
libérés~: la commande TRIM permet au contrôleur d'effacer réellement des
données supprimées, parfois quelques minutes après la suppression. Deux
acquisitions successives d'un SSD allumé peuvent donc différer, sans que
personne ait écrit dessus, et la récupération de fichiers supprimés y est
moins fructueuse que sur un disque magnétique. Les secteurs défectueux d'un
disque magnétique usé, eux, produisent des lectures incomplètes~: l'outil
d'acquisition les journalise et l'enquêteur reporte leur nombre dans son
rapport.

\begin{boxprocedure}
\textbf{Acquérir un disque à froid}
\begin{enumerate}
    \item Photographier la machine, ses câbles et ses étiquettes~; noter le
          numéro de série du disque.
    \item Vérifier le bloqueur d'écriture sur un support de test, et le
          consigner.
    \item Relier le disque au bloqueur, puis à la station d'acquisition.
    \item Choisir le format (E01 par défaut), saisir les métadonnées~:
          numéro de pièce, enquêteur, description.
    \item Acquérir, en calculant SHA-256 (et MD5 si une base l'exige).
    \item Lire le journal de l'outil~: secteurs illisibles, erreurs,
          durée.
    \item Vérifier l'image~: recalculer l'empreinte de l'image et la
          comparer à celle du support.
    \item Faire une copie de travail, vérifier son empreinte, et sceller
          l'image de référence.
\end{enumerate}
\end{boxprocedure}

% ============================================================================
\section{Capturer la Mémoire}
\label{sec:c03-memoire}
% ============================================================================

La mémoire vive contient ce qu'aucun disque ne garde~: la liste des
processus actifs, les connexions réseau ouvertes, les commandes saisies, les
clés de chiffrement des volumes ouverts, et parfois des programmes qui
n'existent qu'en mémoire. On la capture avec un outil exécuté sur la machine
elle-même~: WinPmem ou DumpIt sous Windows, LiME ou AVML sous Linux. Ces
outils sont libres ou gratuits et largement documentés.

\begin{boxinfo}
\textbf{Une acquisition qui ne se répète pas.} La mémoire change pendant la
capture~: les processus continuent de tourner, et l'outil lui-même occupe de
la mémoire. On ne pourra jamais recalculer l'empreinte de la source pour la
comparer à celle de la capture. L'empreinte de la mémoire capturée prouve
seulement que la capture n'a pas changé depuis. C'est pourquoi l'on
documente avec soin l'heure, l'outil, sa version, le support utilisé et ce
qui tournait sur la machine.
\end{boxinfo}

\begin{tableaucours}{L{4.0cm}L{8.2cm}}{%
    \tableauHeader{Précaution} &
    \tableauHeader{Raison}
}
Lancer l'outil depuis un support externe & Ne pas écrire l'outil sur le
disque à examiner \\
Écrire la capture sur un support externe & Ne pas écraser l'espace libre
du disque, où dorment des fichiers supprimés \\
Ne rien installer & Une installation modifie le registre et le disque \\
Noter l'heure système et l'heure réelle & Mesurer la dérive de l'horloge \\
Consigner chaque commande & L'empreinte de l'enquêteur doit pouvoir être
distinguée de celle de l'attaquant \\
\end{tableaucours}
\vspace{-0.8em}
\begin{center}{\small\itshape Tableau~3.4. Précautions de la capture mémoire}\end{center}

% ============================================================================
\section{Les Journaux et les Données des Tiers}
\label{sec:c03-journaux}
% ============================================================================

Les journaux sont souvent la seule trace d'un événement ancien. On les
préserve par un export complet, avec son empreinte, plutôt que par une
recherche filtrée~: un filtre écrit le 14~mars sur les mots « svc\_backup »
aurait laissé de côté les autres comptes, alors qu'on ne savait pas encore
lesquels l'attaquant avait utilisés.

Les données détenues par des tiers (opérateur de paiement mobile, fournisseur
d'accès, hébergeur, éditeur d'un service en nuage) ne se copient pas~: on les
demande, et en matière pénale on les requiert. Le Chapitre~\ref{chap:normes-droit}
précise les fondements de ces demandes. Deux réflexes valent dans tous les
cas~: demander très tôt la conservation des données, avant même leur
communication, et demander comment elles ont été extraites.

% ============================================================================
\section{Les Acquisitions de NYANGA}
\label{sec:c03-nyanga}
% ============================================================================

\begin{boxscenario}{Opération NYANGA. Trois serveurs, trois méthodes}
\textbf{Constatations.}
\begin{itemize}
    \item SRV-AD01~: capture mémoire à chaud le 14~mars (NC-01)~; le serveur,
          indispensable à toutes les connexions, est ensuite remplacé par un
          contrôleur de secours, puis acquis à froid le 16~mars.
    \item SRV-CPT01~: non chiffré~; arrêt contrôlé le 15~mars à 15~h~55,
          acquisition à froid par bloqueur matériel de 16~h~00 à 19~h~25, au
          format E01 (NC-03)~; aucun secteur illisible.
    \item SRV-FIC01~: volume de données chiffré par BitLocker, clé chez le
          prestataire~; capture mémoire puis copie logique du volume
          déverrouillé, à chaud, le 18~mars (NC-07)~; la clé de récupération
          arrive le 22~mars et permet une acquisition à froid complémentaire.
\end{itemize}

\textbf{Analyse} (confiance élevée). Les trois choix sont justifiés par ce
que chaque méthode aurait fait perdre. L'acquisition à chaud de SRV-FIC01 a
modifié la machine (outil exécuté, fichiers temporaires)~; ces
modifications sont consignées heure par heure, ce qui permet de les écarter
de l'analyse. L'acquisition complémentaire du 22~mars confirme que la copie
logique du 18 contenait bien tous les fichiers du volume.
\end{boxscenario}

\begin{boxobjet}{Image du disque de SRV-CPT01, 15 mars 2025}
\textbf{Source~:} serveur SRV-CPT01, disque système de 500~Go.

\textbf{Pièce~:} image E01, scellé NC-03, copie de travail séparée.

\textbf{Méthode~:} arrêt contrôlé, bloqueur matériel vérifié le matin même,
acquisition de 16~h~00 à 19~h~25~; empreinte SHA-256 de la source et de
l'image identiques.

\textbf{Ce qu'elle peut prouver~:} les fichiers présents et supprimés, leurs
dates, les journaux Windows, les traces d'exécution de programmes.

\textbf{Ce qu'elle ne prouve pas~:} ce qui était en mémoire avant l'arrêt du
15~mars, ni ce qui s'est passé sur d'autres serveurs.
\end{boxobjet}

\begin{boxlecture}{Les acquisitions de NYANGA au 22 mars}
\begin{tableaucours}{L{3.4cm}ccc}{%
    \tableauHeader{Question} &
    \tableauHeader{SRV-AD01} & \tableauHeader{SRV-CPT01} &
    \tableauHeader{SRV-FIC01}
}
Mémoire capturée avant arrêt & \OUI{} & \NON{} & \OUI{} \\
Disque acquis sans altération & \OUI{} & \OUI{} & \OUI{} \\
Altérations éventuelles documentées & \OUI{} & \OUI{} & \OUI{} \\
Complétude de l'image vérifiée & \OUI{} & \OUI{} & \OUI{} \\
\end{tableaucours}

\medskip
\textbf{Lecture.} Une seule cote \NON{}~: la mémoire de SRV-CPT01 n'a pas
été capturée avant son arrêt. L'enquête ne pourra donc pas savoir si un
programme de l'attaquant y tournait encore le 15~mars. C'est une limite, et
le rapport la mentionne.
\end{boxlecture}

% ============================================================================
\section{Ce que l'Acquisition demande à chaque Métier}
\label{sec:c03-metiers}
% ============================================================================

\begin{boxmetier}{Analyste en réponse à incident}
Avoir, avant l'incident, une trousse prête~: bloqueur matériel vérifié,
disques neufs, outils de capture mémoire sur un support en lecture seule,
formulaires imprimés. On n'improvise pas une acquisition à trois heures du
matin.
\end{boxmetier}

\begin{boxmetier}{Application OPJ/Expert Judiciaire}
Quand l'acquisition a été faite par d'autres, l'expert ne la refait pas
nécessairement~; il la contrôle~: journal de l'outil, empreintes, bloqueur,
méthode. S'il refait une acquisition, il explique pourquoi.
\end{boxmetier}

\begin{boxmetier}{RSSI}
Conserver, hors des machines, les clés de récupération des volumes chiffrés,
dans un coffre accessible à l'équipe de crise. À NYANGA, quatre jours ont
été perdus à attendre une clé détenue par un prestataire.
\end{boxmetier}

\separateur

\begin{boxresume}
\textbf{Ce qu'il faut retenir de ce chapitre~:}
\begin{itemize}
    \item On recueille de la trace la plus volatile à la moins volatile~;
          les journaux soumis à rotation sont plus fragiles qu'ils ne
          paraissent.
    \item Le choix entre froid et chaud se décide d'après ce que l'on
          perdrait, et se justifie par écrit.
    \item Un disque s'acquiert derrière un bloqueur d'écriture vérifié,
          dans un format qui garde les métadonnées, et l'acquisition se
          vérifie.
    \item Une capture mémoire ne se répète pas~: on documente l'outil,
          l'heure et l'empreinte laissée sur la machine.
    \item À NYANGA, trois serveurs ont reçu trois méthodes différentes,
          chacune justifiée~; la mémoire de SRV-CPT01 manque.
\end{itemize}
\end{boxresume}

\begin{boxinfo}
\textbf{Points de contrôle.} Vous avez compris ce chapitre si vous savez~:
\begin{itemize}
    \item classer des traces selon leur volatilité~;
    \item choisir et justifier une méthode d'acquisition~;
    \item lire le journal d'un outil d'acquisition et y repérer ce qui
          compte.
\end{itemize}
\end{boxinfo}

% ============================================================================
\section*{Exercices}
% ============================================================================

\begin{boxexo}[title={\ding{45}\quad Exercice~3.1. Classer par volatilité \niveaubase}]
Classez de la plus volatile à la moins volatile~: (a) le journal de la
passerelle, conservé 60~jours~; (b) la table ARP du contrôleur de domaine~;
(c) le fichier d'échange de SRV-CPT01~; (d) les sauvegardes sur bande de
février~; (e) les connexions réseau ouvertes de SRV-FIC01~; (f) la capture
de la sonde, conservée 48~heures.
\end{boxexo}

\begin{boxexo}[title={\ding{45}\quad Exercice~3.2. Choisir la méthode \niveaumoyen}]
Pour chaque situation, proposez une méthode d'acquisition et justifiez-la
par ce que l'on perdrait avec l'autre~:
\begin{enumerate}[(a)]
    \item un portable d'agent, éteint, saisi dans un sac~;
    \item le serveur de l'interface de paiement, qui ne peut pas s'arrêter
          plus de dix minutes~;
    \item une machine virtuelle hébergée sur l'hyperviseur du siège~;
    \item la messagerie de l'établissement, hébergée par un fournisseur en
          nuage.
\end{enumerate}
\end{boxexo}

\begin{boxexo}[title={\ding{45}\quad Exercice~3.3. Lire un journal d'acquisition \niveaumoyen}]
Le journal d'une acquisition indique~: « Source~: 500\,107\,862\,016
octets. Secteurs lus avec erreur~: 384. Remplacés par des zéros. Empreinte
SHA-256 de l'image~: 9f3c\ldots{} Vérification~: empreinte de l'image
identique à l'empreinte calculée pendant l'acquisition. »
\begin{enumerate}[(a)]
    \item L'image est-elle intègre~? Est-elle complète~?
    \item Quelle taille de données est perdue, si les secteurs font 512
          octets~?
    \item Comment le rapport doit-il présenter ce fait~?
\end{enumerate}
\end{boxexo}

\begin{boxexo}[title={\ding{45}\quad Exercice~3.4. SRV-FIC01 à chaud \niveauexpert}]
\begin{enumerate}[(a)]
    \item Rédigez la justification écrite du choix d'une acquisition à chaud
          de SRV-FIC01 le 18~mars.
    \item Listez les traces que l'acquisition à chaud a pu laisser sur la
          machine, et dites comment les distinguer de celles de
          l'attaquant.
    \item La clé BitLocker arrive le 22~mars. Que permet l'acquisition à
          froid complémentaire, et quelle comparaison faites-vous entre les
          deux copies~?
\end{enumerate}
\end{boxexo}

\bigskip

\begin{boxinfo}
\textbf{Pour aller plus loin.}
\begin{itemize}
    \item La RFC~3227~\cite{rfc3227} et l'ISO/IEC~27037~\cite{iso27037_2012}.
    \item B.~Carrier, \textit{File System Forensic
          Analysis}~\cite{carrier2005}, chapitre sur l'acquisition.
\end{itemize}
\end{boxinfo}

\begin{boxinfo}
\textbf{Transition.} L'image du disque de SRV-CPT01 est sous scellé, sa
copie de travail est prête. Que contient-elle~? Le Chapitre~\ref{chap:disque}
apprend à lire un système de fichiers et ses artefacts.
\end{boxinfo}
